% ===== PRÉAMBULE CANONIQUE — à coller VERBATIM en tête de CHAQUE .tex =====
\documentclass[11pt,a4paper]{article}
\usepackage[utf8]{inputenc}\usepackage[T1]{fontenc}\usepackage{lmodern}
\usepackage[french]{babel}
\usepackage{amsmath,amssymb,mathtools}\usepackage{icomma}
\usepackage{siunitx}\sisetup{locale=FR}  % \qty{}{}, \si{}, \ang{}
\usepackage{xcolor}\usepackage{colortbl}
\usepackage{tikz}\usepackage{pgfplots}\pgfplotsset{compat=1.18}
\usepackage[siunitx,europeanresistors]{circuitikz} % circuits (élec.)
\usepackage[most]{tcolorbox}\usepackage{enumitem}\usepackage{array,booktabs}
\usepackage[a4paper,margin=2.2cm]{geometry}
\usetikzlibrary{arrows.meta,calc,angles,quotes,positioning,patterns,%
  backgrounds,shapes.geometric,decorations.pathmorphing,decorations.markings,intersections}
\definecolor{primary}{RGB}{222,49,99}\definecolor{primarydark}{RGB}{150,22,55}
\definecolor{defcol}{RGB}{0,114,178}\definecolor{methcol}{RGB}{52,92,148}
\definecolor{excol}{RGB}{0,135,90}\definecolor{attncol}{RGB}{200,40,55}
\tcbset{boxbase/.style={enhanced,breakable,boxrule=0.9pt,arc=2.5pt,
  left=3mm,right=3mm,top=2.3mm,bottom=2.3mm,fonttitle=\bfseries,coltitle=white}}
% --- boîtes TUTORIEL (noms EXACTS, ne pas renommer) ---
\newtcolorbox{cle}[1][]{boxbase,colback=defcol!6,colframe=defcol,
  title={\textsc{À retenir}\ifx\relax#1\relax\relax\else\textmd{ : \itshape #1}\fi}}
\newtcolorbox{methode}[1][]{boxbase,colback=methcol!7,colframe=methcol,sharp corners,
  title={\textsc{Méthode}\ifx\relax#1\relax\relax\else\textmd{ --- \itshape #1}\fi}}
\newtcolorbox{exemplebox}[1][]{boxbase,colback=excol!5,colframe=excol,
  title={\textsc{Exemple}\ifx\relax#1\relax\relax\else\textmd{ : \itshape #1}\fi}}
\newtcolorbox{piege}[1][]{boxbase,colback=attncol!6,colframe=attncol,
  title={\textsc{Piège}\ifx\relax#1\relax\relax\else\textmd{ : \itshape #1}\fi}}
\newtcolorbox{remarque}[1][]{boxbase,colback=black!4,colframe=black!45,
  title={\textsc{Remarque}\ifx\relax#1\relax\relax\else\textmd{ : \itshape #1}\fi}}
% --- boîtes EXERCICE / CORRIGÉ (noms EXACTS, auto-comptées) ---
\newtcolorbox[auto counter]{exo}[1][]{boxbase,colback=primary!4,colframe=primary,
  title={\textsc{Exercice}~\thetcbcounter\ifx\relax#1\relax\relax\else\textmd{ --- \itshape #1}\fi}}
\newtcolorbox[auto counter]{corrige}[1][]{boxbase,colback=excol!4,colframe=excol,
  title={\textsc{Corrigé}~\thetcbcounter\ifx\relax#1\relax\relax\else\textmd{ --- \itshape #1}\fi}}
\newcommand{\vv}[1]{\vec{#1}}\newcommand{\ds}{\displaystyle}
\newcommand{\R}{\mathbb{R}}\newcommand{\N}{\mathbb{N}}\newcommand{\Z}{\mathbb{Z}}
% (un \usepackage{fancyhdr}+en-tête est possible ; il est ignoré côté web)
% =========================================================================
\begin{document}
\sisetup{per-mode=symbol}

\begin{exo}[calcul direct de l'énergie cinétique]
Un ballon de masse $m=\qty{2}{\kilogram}$ se déplace à la vitesse $v=\qty{3}{\metre\per\second}$.
Calcule son énergie cinétique $E_c$.
\end{exo}

\begin{exo}[énergie cinétique avec conversion de vitesse]
Un cycliste et son vélo ont une masse totale $m=\qty{80}{\kilogram}$ et roulent à
$v=\qty{18}{\kilo\metre\per\hour}$.
\begin{enumerate}[label=\alph*)]
  \item Convertis la vitesse en \si{\metre\per\second}.
  \item Calcule l'énergie cinétique de l'ensemble.
\end{enumerate}
\end{exo}

\begin{exo}[l'effet du carré de la vitesse]
À la vitesse $v$, un objet possède une énergie cinétique $E_c=\qty{20}{\joule}$.
\begin{enumerate}[label=\alph*)]
  \item Quelle est son énergie cinétique s'il roule à la vitesse $2v$ ?
  \item Quelle est son énergie cinétique s'il roule à la vitesse $3v$ ?
\end{enumerate}
\end{exo}

\begin{exo}[théorème de l'énergie cinétique à partir du repos]
Un palet de masse $m=\qty{5}{\kilogram}$, initialement \emph{au repos}, reçoit un travail total
$\sum W=\qty{250}{\joule}$.
\begin{enumerate}[label=\alph*)]
  \item Que vaut l'énergie cinétique finale $E_c$ du palet ?
  \item Déduis-en sa vitesse finale $v$.
\end{enumerate}
\end{exo}

\begin{exo}[retrouver la vitesse à partir de l'énergie cinétique]
Une bille de masse $m=\qty{5}{\kilogram}$ possède une énergie cinétique $E_c=\qty{90}{\joule}$.
En partant de $E_c=\tfrac12 m v^2$, calcule sa vitesse $v$.
\end{exo}
\end{document}
